\section{Introduction}

Bayesian component separation in the pixel domain fits, in every pixel, the amplitude of each sky
component, and for the whole sky or for large regions a small number of spectral parameters that say
how each component scales with frequency. In the Gibbs samplers that the Planck, BeyondPlanck and
Cosmoglobe analyses are built on \citep{eriksen2008,svalheim2023,watts2023}, the two kinds of
parameter are updated in turn. The amplitudes are many and each is poorly determined; the spectral
indices are few and are what the cleaning of large-scale polarization, and therefore the optical
depth to reionization, depends on. The published polarized synchrotron indices from the same public
maps span $-2.9$ to $-3.3$ depending on region, band set and method
\citep{dunkley2009,fuskeland2014,svalheim2023,weiland2022}, which is larger than their quoted errors.

This paper is about the statistical step that sits between the maps and such an index: how the
amplitudes are removed. That step has three places where it can go wrong without any error message,
and we treat them in turn.

\emph{The sampler} (Sect.~\ref{sec:legacy}). A move that proposes a new index, re-solves the
amplitudes and accepts on the $\chi^2$ of the new pair looks like a way to move both together. We
show that on its own it is a valid pseudo-marginal kernel for the profile posterior, that the
amplitude step of the Gibbs sweep turns it into Monte Carlo within Metropolis, and that the chain
then converges to the profile posterior raised to the power $1/T$ with
$T = \sqrt{\pi k/8} + \tfrac12$ for $k$ amplitudes. The index error bar from such a chain is too wide
by $0.79\,k^{1/4}$, a factor 16 for the $1.9\times10^5$ amplitudes of an $N_{\rm side}=64$ analysis,
with an acceptance rate that stays near one half.

\emph{The prior} (Sect.~\ref{sec:flat}). With flat priors on the amplitudes, the standard choice, the
index posterior contains the volume of amplitude space. This is known to bias indices
\citep{eriksen2008,dunkley2009,alonso2017}. We show that it is exactly equivalent to an exponential
prior on the index whose rate is the number of amplitudes times the logarithm of the reference
frequency at which the amplitudes are defined, give the resulting shift in closed form, and show that
its sign flips between the two reference frequencies in common use, 23 and 30 GHz, for the WMAP and
LFI bands when the CMB amplitude is constrained (with a free CMB amplitude both pull the index flat).
The shift does not shrink as pixels are added.

\emph{The error bar} (Sect.~\ref{sec:profile}). The profile likelihood has neither problem and its
score is exactly unbiased, but its curvature understates the error where the amplitudes are faint.
We give the correction in closed form, show that the obvious way to evaluate it on real data is
biased in the wrong direction, give the version that is not, and show its coverage over four decades
of the controlling ratio. A zero-mean amplitude prior with a free scale turns out to be the same
estimator when the noise pattern is uniform, and a biased one when it is not.

Sect.~\ref{sec:data} applies the estimators side by side to the WMAP nine-year $K$ to $V$ band
polarization maps with their released dense noise covariances, on regions defined by geometry with
the brightest emission masked, so that the difference between their answers is the difference between
the estimators. It then lists what else moves the index on these maps and by how much, where the
model with one power law per region passes and where it fails, and why we give the indices as the
answer of that model and not as a measurement. Sect.~\ref{sec:limits} states where the
recommended estimator stops being enough, and Sect.~\ref{sec:impl} describes the implementation and
its tests. All results are for maps-only analyses at low resolution with amplitudes that enter
linearly; we do not treat the time-ordered data.
